\section{Objetivos}
\label{sec:objetivos}

\subsection{Objetivo general}

Desarrollar un modelo de predicción de heladas mediante aprendizaje automático con series temporales, para anticipar su ocurrencia a 24 y 48 horas en cultivos altoandinos de Cusco, considerando recall, precisión, F1, AUC-PR y MCC frente a un modelo base de persistencia.
\subsection{Objetivos específicos}

\begin{itemize}
  \item \textbf{OE1.} Recopilar y preprocesar los datos históricos de estaciones meteorológicas altoandinas de la región Cusco.
  \item \textbf{OE2.} Identificar las variables meteorológicas y características temporales con mayor poder predictivo de heladas.
  \item \textbf{OE3.} Entrenar y comparar modelos de aprendizaje automático (regresión logística, Random Forest, XGBoost y LSTM) para determinar el de mejor desempeño.
  \item \textbf{OE4.} Evaluar el desempeño del modelo a 24 y 48 horas de anticipación frente a un modelo base de persistencia, usando recall, precisión, F1, AUC-PR y MCC.
\end{itemize}

\subsection{Correspondencia entre problemas y objetivos}

\begin{center}
\begin{tabular}{p{0.45\textwidth}p{0.45\textwidth}}
\toprule
\textbf{Problema} & \textbf{Objetivo que lo responde} \\
\midrule
Problema general & Objetivo general \\
PE1: calidad de los datos & OE1: recopilar y preprocesar datos \\
PE2: variables relevantes & OE2: identificar variables predictivas \\
PE3: mejor algoritmo & OE3: entrenar y comparar modelos \\
PE4: desempeño a 24 y 48 h & OE4: evaluar frente al modelo base \\
\bottomrule
\end{tabular}
\end{center}